% Editar este fichero y rellenar con el resumen del TFM

Este Trabajo de Fin de Máster presenta el diseño, desarrollo y evaluación de un sistema 
orientado a la automatización de la documentación médica mediante la 
transformación de dictados de voz en informes clínicos normalizados. 
El sistema propone un flujo de trabajo \englishText{pipeline} que enumera las 
siguientes etapas:
\begin{enumerate}
    \setlength{\itemsep}{1pt}
    \setlength{\parsep}{0pt}
    \setlength{\topsep}{4pt}
    \item \textbf{Captura segura de audio:} 
        Implementación de una arquitectura cliente-servidor con cifrado de extremo a extremo
        para garantizar la integridad de los datos.
    \item \textbf{Reconocimiento automático del habla (\englishText{ASR}, Automatic Speech Recognition):} 
        Transcripción del audio mediante modelos de inteligencia artificial
        adaptados al léxico médico.
    \item \textbf{Extracción y normalización de entidades clínicas:}
        Identificación, categorización y normalización de conceptos clínicos bajo estándares
        internacionales como \englishText{SNOMED CT} (\englishText{Systematized Nomenclature of Medicine -- Clinical Terms})
        y la \englishText{CIE-10-ES} (Clasificación Internacional de Enfermedades, 10.\textsuperscript{a} revisión,
        modificación clínica española).
    \item \textbf{Generación de informes mediante \englishText{LLMs} (Large Language Models):}
        Síntesis de información utilizando modelos de lenguaje de gran escala a partir de un
        contexto verificado de entidades clínicas previamente normalizadas, con el fin de
        reducir el riesgo de invención de información ausente en el dictado original.
    \item \textbf{Salida estructurada e interoperable:}
        Exportación de resultados en formatos estándares (\englishText{JSON}/\englishText{FHIR},
        \englishText{Fast Healthcare Interoperability Resources})
        para su integración directa en sistemas de Historia Clínica Electrónica (\englishText{HCE}) y
        generación de informes clínicos.
\end{enumerate}

El proyecto se fundamenta en los principios de privacidad desde el diseño,
asegurando el cumplimiento riguroso de los marcos legales vigentes (\englishText{RGPD}, Reglamento General de
Protección de Datos, y \englishText{LOPDGDD}, Ley Orgánica de Protección de Datos Personales y garantía de los
derechos digitales).
La viabilidad de la propuesta se valida a través de una evaluación cuantitativa basada en métricas de precisión técnica
(\englishText{WER}, \englishText{Word Error Rate}, para la transcripción y puntuación F$_{1}$ para la extracción de entidades),
así como indicadores de eficiencia temporal y relevancia clínica.